\subsection{Indirect Support Predicts Direct Genetic Support}

\begin{figure}[t!]
    \centering
    {{canva:4}}
    \caption{Validation of indirect genetic support. (a) Leave-one-chromosome-out (LOCO) cross-validation across 12 common traits: indirect support learned from all other chromosomes predicts direct support on the held-out chromosome. Dashed line indicates $y = x$. (b) Rare disease validation: cumulative match characteristic (CMC) curves for recovery of held-out OMIM disease genes across 40 rare diseases. (c) Temporal validation: improvement in model fit ($\Delta R^2$) for post-2018 MAGMA Z-scores as a function of indirect support threshold, using pre-2018 annotations. (d) Distribution of the number of significant genesets per trait. (e) Semantic similarity between significant genesets and their associated phenotype versus a random null.}
    \label{fig:validation-genetics}
\end{figure}

% Panel A and B: indirect predicts direct in common and rare disease
We first evaluated whether indirect support predicts direct genetic support across both common and rare traits. For common traits, we tested this using leave-one-chromosome-out (LOCO) cross-validation across 12 validation traits (Supplementary Table \ref{tab:common-validation-traits}). In each fold, indirect support was learned from all chromosomes except one, then used to predict direct genetic support on the held-out chromosome (Figure \ref{fig:validation-genetics}a). Indirect support significantly predicted direct support for every held-out chromosome across all 12 traits (Supplementary Table \ref{tab:loco-validation}; trait-level $\beta$ range {{MIN_LOCO_BETA}} to {{MAX_LOCO_BETA}}, all $p < {{MAX_LOCO_BETA_P}}$). The fitted relationships consistently fell below the $y = x$ calibration line, reflecting the expected conservative bias: indirect support estimates gene-trait relevance, which is broader than direct association, so not all indirectly supported genes will carry detectable associations at current sample sizes.

For rare diseases, we performed an analogous evaluation across 40 rare traits (Supplementary Table \ref{tab:rare-validation-traits}) using OMIM disease-gene annotations curated by Orphanet as a proxy for direct support. For each disease, we held out one known disease gene and asked how often it was recovered among the top-$K$ indirectly prioritized genes, summarized by a cumulative match characteristic (CMC) curve (Figure \ref{fig:validation-genetics}b). Indirect support recovered held-out rare disease genes with a CMC-AUC of {{RARE_CMC_AUC}} (95\% CI: {{RARE_CMC_AUC_CI_LOWER}} to {{RARE_CMC_AUC_CI_UPPER}}) and a Mean Reciprocal Rank (MRR) of {{RARE_MRR}} (95\% CI: {{RARE_MRR_CI_LOWER}} to {{RARE_MRR_CI_UPPER}}), implying the true gene is typically ranked within the top 5 to 10 candidates.

% Panel C: temporal validation
Next, we sought to control for potential circularity caused by genetic annotations that may have been informed by genetic data. We curated a temporal validation set of 67 common traits that had a GWAS published both before and after 2018, with the post-2018 GWAS having higher statistical power. We applied the Mouse + MSigDB model using pre-2018 versions of both annotation libraries. We then calculated gene-level MAGMA Z-scores from both the pre- and post-2018 GWAS for each trait and regressed the post-2018 MAGMA Z-scores on the pre-2018 scores while progressively restricting to genes meeting increasing thresholds of indirect support (Figure \ref{fig:validation-genetics}c). The improvement in model fit ($\Delta R^2$) increased significantly with the indirect support threshold ($\beta$ = {{TEMPORAL_DELTA_R2_BETA}}, $p$ = {{TEMPORAL_DELTA_R2_BETA_P}}), demonstrating that indirect support is predictive of future GWAS discoveries and is not simply recapitulating latent genetic information embedded in curated libraries.

% Panel D and E: interpretability
We then assessed the interpretability of PIGEAN's inferred geneset effects, which govern the calculation of indirect support. Across {{NUM_TRAITS_MOUSE_MSIGDB_CONVERGED}} converged traits, approximately {{PERCENTAGE_OF_TRAITS_WITH_LESS_THAN_100_GENESETS}}\% ({{NUM_TRAITS_WITH_LESS_THAN_100_GENESETS}} of {{NUM_TRAITS_MOUSE_MSIGDB_CONVERGED}}) had fewer than 100 genesets with significant joint effect sizes ($\beta > 0.01$; Figure~\ref{fig:validation-genetics}d), indicating that PIGEAN produces sparse and interpretable outputs. The significant genesets were semantically related to their associated phenotypes: the mean cosine similarity between significant geneset names and phenotype descriptions was {{MEAN_PHENOTYPIC_SIMILARITY_OF_SIGNIFICANT_GENESETS}} compared to {{MEAN_PHENOTYPIC_SIMILARITY_OF_RANDOM_NULL}} for a random null ($t$ = {{T_STATISTIC_PHENOTYPIC_SIMILARITY}}, $p \ll 0.001$; Figure~\ref{fig:validation-genetics}e).
